
        %%%%% ~~~~~~~~~~~~~~~~~~~~ %%%%%

\section{From generalized totals defined at level of estimation units to those defined at level of imputation units}
\label{GEstToOneStageClusterSampling}

        %%%%% ~~~~~~~~~~~~~~~~~~~~ %%%%%

\subsection{Generalized totals defined at the level of estimation units}
\label{Whatever}
\setcounter{theorem}{0}

        %%%%% ~~~~~~~~~~~~~~~~~~~~ %%%%%

Statistics Canada's Integrated Business Survey Program (IBSP)
uses a three-level hierarchy of ``units'':
\begin{itemize}
\item
    sampling units
    \begin{itemize}
    \item
        The sampling design is defined at the level of sampling units.
    \item
        Each sampling unit is composed of one or more imputation units.
    \end{itemize}
\item
    imputation units
    \begin{itemize}
    \item
        Collection and imputation is performed at the level of imputation units.
    \item
        Each imputation unit is composed of one or more estimation units.
    \item
        It is assumed that the sampling design is {\color{red}one-stage cluster sampling},
        in the sense that, the sample is a probability sample of \textit{sampling units}, and
        for each selected \textit{sampling unit},
        all of its constituent \textit{imputation units} are selected.
        In standard cluster sampling terminology,
        the \textit{sampling units} are called the \textit{primary sampling units} (PSU),
        while the \textit{imputation units} are called the \textit{secondary sampling units} (SSU).
    \end{itemize}
\item
    estimation units
    \begin{itemize}
    \item
        In practice, {\color{red}domain membership indicators} are defined
        at the level of estimation units.
    \end{itemize}
\end{itemize}
G-Est is a Statistics Canada generalized system that computes (among other functionalities)
point and variance estimates (from probability sample data) for population parameters of the following form:

        %%%%% ~~~~~~~~~~~~~~~~~~~~ %%%%%

\vskip 0.5cm
\begin{definition}[Generalized total defined at the level of estimation units]
\mbox{}
\vskip 0.1cm
\noindent
A \textbf{generalized total defined at the level of estimation units}
is a population parameter of the following form:
\begin{equation*}
T_{\alpha Y}
\;\; := \;\;
    \underset{i \in U}{\sum}\;\;
    \underset{j \in \mathcal{C}_{i}}{\sum}\;
    \underset{k \in \mathcal{E}_{ij}}{\sum}\;
    \alpha_{ijk}Y_{ijk}\,,
\end{equation*}
where
\begin{itemize}
\item
    $U$\, is the population of sampling units (i.e., primary sampling units),
\item
    $\mathcal{C}_{i}$,\, for \,$i \in U$,\,
    is the collection of imputation units (i.e., secondary sampling units)
    belonging to sampling unit $i \in U$,
\item
    $\mathcal{E}_{ij}$,\, for \,$i \in U$\, and \,$j \in \mathcal{C}_{i}$,\,
    is the collection of estimation units
    belonging to imputation unit \,$j \in \mathcal{C}_{i}$\, of sampling unit $i \in U$,
\item
    $Y_{ijk} \in \Re$\, is the value
    of an ($\Re$-valued) population characteristic \,$\mathbf{Y}$\,
    defined at the level of the estimation units; more precisely,
    \,$Y_{ijk}$\, is the value of \,$\mathbf{Y}$\,
    for estimation unit \,$k \in \mathcal{E}_{ij}$\,
    of imputation unit \,$j \in \mathcal{C}_{i}$\,
    belonging to sampling unit \,$i \in U$,\,
\item
    $\alpha_{ijk} \in \Re$\, is the coefficient defined at the level of the estimation unit.
\end{itemize}
\end{definition}

        %%%%% ~~~~~~~~~~~~~~~~~~~~ %%%%%

\vskip 1.0cm
\subsection{Rewriting generalized totals defined at level of estimation units as sums over imputation units via allocation factors}
\label{Whatever}
\setcounter{theorem}{0}

        %%%%% ~~~~~~~~~~~~~~~~~~~~ %%%%%

\vskip 0.2cm
\begin{lemma}
\label{GeneralizedTotalsAsSumsOverImputationUnits}
\mbox{}
\vskip 0.1cm
\noindent
Let
\begin{equation*}
T_{\alpha Y}
\;\; := \;\;
    \underset{i \in U}{\sum}\;\;
    \underset{j \in \mathcal{C}_{i}}{\sum}\;
    \underset{k \in \mathcal{E}_{ij}}{\sum}\;
    \alpha_{ijk}Y_{ijk}\,,
\end{equation*}
be a generalized total defined at the level of estimation units.
\vskip 0.1cm
\noindent
Suppose that there exists a collection
\begin{equation*}
\left\{\;\,
    f_{ijk} \in \Re
    \;\,\left\vert\,
    \begin{array}{c} i \in U, \\ j \in \mathcal{C}_{i}, \\ k \in \mathcal{E}_{ij} \end{array}
    \right.
    \!\!\right\}
\end{equation*}
of real numbers such that
\begin{equation*}
Y_{ijk}
\; = \;
    f_{ijk}\cdot\!\left(\,
        \underset{l\in\mathcal{E}_{ij}}{\sum}\,Y_{ijl}
        \right),
\quad
\textnormal{for each \,$i \in U$,\, \,$j \in \mathcal{C}_{i}$,\, and \,$k \in \mathcal{E}_{ij}$}
\end{equation*}
Then, the generalized total \,$T_{\alpha Y}$\, (defined at the level of estimation units)
can be re-expressed as a sum over the imputation units; more precisely,
\begin{equation*}
T_{\alpha Y}
\;\; := \;\;
    \underset{i \in U}{\sum}\;\;
    \underset{j \in \mathcal{C}_{i}}{\sum}\;
    \underset{k \in \mathcal{E}_{ij}}{\sum}\;
    \alpha_{ijk}Y_{ijk}
\;\; = \;\;
    \underset{i \in U}{\sum}\;
    \underset{j \in \mathcal{C}_{i}}{\sum}\;
    \alpha_{ij}Y_{ij}\,,
\end{equation*}
where
\begin{equation*}
\alpha_{ij}
\;\, := \;
    \underset{k \in \mathcal{E}_{ij}}{\sum}\,\alpha_{ijk}\,f_{ijk}
\quad\quad
\textnormal{and}
\quad\quad
Y_{ij}
\;\, := \;
    \underset{k \in \mathcal{E}_{ij}}{\sum}Y_{ijk}
\end{equation*}
\end{lemma}
\proof
% \begin{eqnarray*}
% T_{\alpha Y}
% & := &
%     \underset{i \in U}{\sum}\;\;
%     \underset{j \in \mathcal{C}_{i}}{\sum}\;
%     \underset{k \in \mathcal{E}_{ij}}{\sum}\;
%     \alpha_{ijk}Y_{ijk}
% \;\; = \;\;
%     \underset{i \in U}{\sum}\;
%     (\alpha Y)_{i}
% \\
% Y_{ijk}
% & = &
%     f^{(Y)}_{ijk}\,Y_{ij}\,,
%     \quad
%     \textnormal{where}\;\;
%     \underset{k \in \mathcal{E}_{ij}}{\sum}f^{(Y)}_{ijk} \,=\, 1
% \\
% Y_{ij}
% & = &
%     \underset{k \in \mathcal{E}_{ij}}{\sum}Y_{ijk}
% \;\; = \;\;
%     \underset{k \in \mathcal{E}_{ij}}{\sum}f^{(Y)}_{ijk}\,Y_{ij}
% \;\; = \;\;
%     \left(\,\underset{k \in \mathcal{E}_{ij}}{\sum}f^{(Y)}_{ijk}\right)
%     \!\cdot
%     Y_{ij}
% \;\; = \;\;
%     Y_{ij}
% \end{eqnarray*}
% Hence,
Straightforward algebra:
\begin{eqnarray*}
T_{\alpha Y}
& := &
    \underset{i \in U}{\sum}\;\;
    \underset{j \in \mathcal{C}_{i}}{\sum}\;
    \underset{k \in \mathcal{E}_{ij}}{\sum}\;
    \alpha_{ijk}Y_{ijk}
\;\; = \;\;
    \underset{i \in U}{\sum}\;\;
    \underset{j \in \mathcal{C}_{i}}{\sum}\;
    \underset{k \in \mathcal{E}_{ij}}{\sum}\;
    \alpha_{ijk}\,\left(f_{ijk}\,Y_{ij}\right)
\;\; = \;\;
    \underset{i \in U}{\sum}\;
    \underset{j \in \mathcal{C}_{i}}{\sum}\!
    \left(\,
        \underset{k \in \mathcal{E}_{ij}}{\sum}\,
        \alpha_{ijk}\,f_{ijk}
        \right)
    \!\cdot Y_{ij}
\\
& = &
    \underset{i \in U}{\sum}\;
    \underset{j \in \mathcal{C}_{i}}{\sum}\;
    \alpha_{ij}Y_{ij}
\end{eqnarray*}
\vskip -0.5cm
\qed

        %%%%% ~~~~~~~~~~~~~~~~~~~~ %%%%%

\begin{remark}
\mbox{}
\vskip 0.1cm
\noindent
\begin{itemize}
\item
    The \,$f_{ijk} \in \Re$\, are known as \textbf{allocation factors}.
\item
    For each \,$i \in U$\, and \,$j \in \mathcal{C}_{i}$,\, we have:
    \begin{equation*}
    \underset{k\in\mathcal{E}_{ij}}{\sum}Y_{ijk}
    \; =: \;
        Y_{ij}
    \; \neq \;
        0
    \quad\Longrightarrow\quad
    \textnormal{
        $f_{ijk}$\,
        is uniquely determined as
        \,$
        f_{ijk}
        \,=\,
            \dfrac{Y_{ijk}}{
                \overset{{\color{white}.}}{Y_{ij}}
                }
        \,=\,
            \dfrac{Y_{ijk}}{
                \overset{{\color{white}.}}{\underset{k\in\mathcal{E}_{ij}}{\sum}Y_{ijk}}
                }
        $
        }
    % \underset{k\in\mathcal{E}_{ij}}{\sum}f_{ijk}
    % \; = \;
    %     1
    \end{equation*}
    \vskip 0.1cm
    \noindent
    In particular,
    \;\;$
    \underset{k\in\mathcal{E}_{ij}}{\sum}Y_{ijk}
    \; =: \;
        Y_{ij}
    \; \neq \;
        0
    \;\;\Longrightarrow\;\;
    \underset{k\in\mathcal{E}_{ij}}{\sum}f_{ijk}
    \; = \;
        1
    $,\,
\item
    By Lemma \ref{GeneralizedTotalsAsSumsOverImputationUnits},
    we may regard ageneralized total \,$T_{\alpha Y}$
    (\textit{a priori} defined as a sum over the estimation units)
    as a sum over the imputation units.
    This allows us to simplify the implementation of SEVANI
    by suppressing the estimation units, i.e.,
    by considering only a two-level hierarchy of units:
    the sampling units and the imputation units.
\end{itemize}
\end{remark}

        %%%%% ~~~~~~~~~~~~~~~~~~~~ %%%%%
